
\documentclass{beamer}

\usepackage{xcolor}
\usepackage{color}
\usepackage{graphicx}
\def\xcolorversion{2.00}
\def\xkeyvalversion{1.8}
\setbeamersize{text margin left=0pt,text margin right=0pt}

\makeatletter
\newcommand*\bigcdot{\mathpalette\bigcdot@{.9}}
\newcommand*\bigcdot@[2]{\mathbin{\vcenter{\hbox{\scalebox{#2}{$\m@th#1\bullet$}}}}}
\makeatother

\usepackage{TIKZ_Macra}
\usepackage{xparse}
\usepackage[utf8]{inputenc}
\usepackage[T1]{fontenc}

\useinnertheme[shadow=true]{rounded}
\setbeamercolor{block title}{use=structure,fg=structure.fg,bg=structure.fg!30!bg}
\setbeamercolor{block body}{parent=normal text,use=block title,bg=block title.bg!50!bg}
\setbeamercolor{block title example}{use=example text,fg=example text.fg,bg=example text.fg!30!bg}
\setbeamercolor{block body example}{parent=normal text,use=block title example,bg=block title example.bg!50!bg}

\usepackage[most]{tcolorbox}
\usepackage{lmodern}
\usepackage{mathtools}
\usepackage{amsfonts}
\usepackage{amssymb}
\usepackage{latexsym}
\usepackage{hyperref}
\usepackage{changepage}

\newtcbox{\mybox}{blank, on line, opacitytext=0.5}

\usepackage{Moje_makra}
\usepackage{Moje_Zmienne}
\usepackage{Zbiory_i_konfiguracje}
\usepackage{Complexity}
\usepackage{Sieci}

\newcommand{\Tr}{\mathsf{Tr}}
\newcommand{\Runs}{\mathsf{Runs}}
\newcommand{\PreA}{\mathsf{Pre}}
\newcommand{\PostA}{\mathsf{Post}}
\newcommand{\Part}{\mathcal{P}}
\newcommand{\Blocks}{\mathsf{Blocks}}
\newcommand{\Sat}{\mathsf{Sat}}
\newcommand{\true}{\mathsf{true}}
\newcommand{\false}{\mathsf{false}}
\newcommand{\Act}{\Sigma}
\newcommand{\AP}{\AProp}
\newcommand{\impl}{\mathsf{Impl}}
\newcommand{\absys}{\mathsf{Abs}}
\newcommand{\depth}{\mathsf{md}}

\begin{document}

\begin{frame}[fragile]{}
\Margines{
\begin{center}
  {\huge Theory of Concurrency}\medskip

  {\large Lecture 3c: HML and bisimulation}\medskip

  Definitions and invariance of CTL/CTL*\medskip

  Piotr Hofman
\end{center}
}
\end{frame}


\begin{frame}{Labelled transition systems}
\Margines{
\begin{block}{Definition}
A labelled transition system is
\[
  T=(S,I,\Sigma,\to),
\]
where \(S\) is a set of states, \(I\subseteq S\), \(\Sigma\) is the alphabet of actions, and
\[
  \to\subseteq S\times\Sigma\times S.
\]
We write \(s\xrightarrow{a}t\).
\end{block}

\begin{block}{Question}
When do two states have the same branching behaviour?
\end{block}
}
\end{frame}

\begin{frame}{Bisimulation}
\Margines{
\begin{block}{Definition}
A relation \(B\subseteq S\times S\) is a bisimulation if, whenever \((s,t)\in B\), then for every \(a\in\Sigma\):
\begin{enumerate}
  \item if \(s\xrightarrow{a}s'\), then there exists \(t\xrightarrow{a}t'\) such that \((s',t')\in B\);
  \item if \(t\xrightarrow{a}t'\), then there exists \(s\xrightarrow{a}s'\) such that \((s',t')\in B\).
\end{enumerate}
\end{block}

\begin{block}{Bisimilarity}
We write \(s\sim t\) if there exists a bisimulation \(B\) with \((s,t)\in B\).
\end{block}
}
\end{frame}

\begin{frame}{Hennessy--Milner logic}
\Margines{
\begin{block}{Syntax}
HML formulas are generated by
\[
\varphi ::= \top \mid \neg\varphi \mid \varphi_1\land\varphi_2 \mid \langle a\rangle\varphi,
\qquad a\in\Sigma.
\]
\end{block}

\begin{block}{Derived box}
\[
  [a]\varphi \equiv \neg\langle a\rangle\neg\varphi.
\]
\end{block}

\begin{block}{Reading}
\begin{itemize}
  \item \(\langle a\rangle\varphi\): there exists an \(a\)-successor satisfying \(\varphi\).
  \item \([a]\varphi\): every \(a\)-successor satisfies \(\varphi\).
\end{itemize}
\end{block}
}
\end{frame}

\begin{frame}{Semantics of HML}
\Margines{
\begin{align*}
 s&\models \top,\\
 s&\models \neg\varphi &&\text{iff } s\not\models\varphi,\\
 s&\models \varphi\land\psi &&\text{iff }s\models\varphi\text{ and }s\models\psi,\\
 s&\models \langle a\rangle\varphi
 &&\text{iff there exists }s'\text{ such that }s\xrightarrow{a}s'\text{ and }s'\models\varphi.
\end{align*}

\begin{block}{Example}
\[
  \langle a\rangle(\langle b\rangle\top\land \langle c\rangle\top)
\]
means: after one \(a\)-move there is a state from which both \(b\) and \(c\) are possible.
\end{block}
}
\end{frame}

\begin{frame}{HML is invariant under bisimulation}
\Margines{
\begin{theorem}
If \(s\sim t\), then for every HML formula \(\varphi\),
\[
  s\models\varphi
  \quad\Longleftrightarrow\quad
  t\models\varphi.
\]
\end{theorem}

\begin{block}{Proof idea}
By structural induction on \(\varphi\). The only non-trivial case is \(\langle a\rangle\varphi\), where the matching transition is provided by the bisimulation condition.
\end{block}
}
\end{frame}

\begin{frame}{Hennessy--Milner theorem}
\Margines{
\begin{theorem}
For image-finite LTSs, and in particular for finite LTSs,
\[
  s\sim t
  \quad\Longleftrightarrow\quad
  \forall\varphi\in HML:\;
  s\models\varphi \iff t\models\varphi.
\]
\end{theorem}

\begin{block}{Interpretation}
Bisimulation is exactly the equivalence induced by HML observations.
\end{block}

\begin{block}{Modal depth}
Formulas of modal depth \(k\) observe branching up to depth \(k\).
\end{block}
}
\end{frame}

\begin{frame}{Bisimulation for Kripke structures}
\Margines{
\begin{block}{Definition}
Let \(M=(S,I,R,L)\) be a Kripke structure. A relation \(B\subseteq S\times S\) is a bisimulation if \((s,t)\in B\) implies:
\begin{enumerate}
  \item \(L(s)=L(t)\),
  \item for every \(s'\) with \((s,s')\in R\), there exists \(t'\) with \((t,t')\in R\) and \((s',t')\in B\),
  \item symmetrically, for every \(t'\) with \((t,t')\in R\), there exists \(s'\) with \((s,s')\in R\) and \((s',t')\in B\).
\end{enumerate}
\end{block}
}
\end{frame}

\begin{frame}{CTL and CTL* are bisimulation invariant}
\Margines{
\begin{theorem}
Let \(s,t\) be states of a Kripke structure. If \(s\sim t\), then for every CTL* state formula \(\varphi\),
\[
  M,s\models\varphi
  \quad\Longleftrightarrow\quad
  M,t\models\varphi.
\]
\end{theorem}

\begin{block}{Consequences}
\begin{itemize}
  \item The same holds for every CTL formula.
  \item Bisimilar states cannot be distinguished by branching-time temporal logic.
  \item The bisimulation quotient preserves CTL and CTL* properties.
\end{itemize}
\end{block}
}
\end{frame}

\begin{frame}{Why this is true}
\Margines{
\begin{block}{Main idea}
Bisimulation matches the whole unfolding tree of possible futures.
\end{block}

\begin{itemize}
  \item Atomic propositions are preserved by \(L(s)=L(t)\).
  \item Boolean connectives are handled by induction.
  \item The path quantifiers \(A\) and \(E\) are preserved because every successor choice can be matched on the other side.
  \item Temporal operators are evaluated along matched paths.
\end{itemize}

\begin{block}{Conclusion}
CTL and CTL* can inspect branching, but not beyond bisimulation.
\end{block}
}
\end{frame}

\begin{frame}{Summary}
\Margines{
\begin{itemize}
  \item HML is a small modal logic for one-step branching observations.
  \item Bisimulation preserves all HML formulas.
  \item On finite systems, HML equivalence coincides with bisimulation.
  \item CTL and CTL* are also invariant under bisimulation.
\end{itemize}
}
\end{frame}

\end{document}
